\section{Zadanie 12}

Kasyno testuje nową maszynę. Aby zagrać gracz płaci \$\(1\), zaś maszyna z prawdopodobieństwem \(\frac{4}{25}\) wypłaca \$\(3\), a prawdopodobieństwem \(\frac{1}{200}\) wypłaca \$\(100\), a w pozostałych przypadkach nic nie wypłaca. Ku zaskoczeniu pracowników kasyna w ciągu pierwszego miliona gier maszyna przegrała \$\(10000\). Oszacuj przy pomocy nierówności Czernowa prawdopodobieństwo zajścia takiego zdarzenia. Numerycznie wyznacz najlepsze możliwe ograniczenie (znajdź odpowiednie \(t\)).

\subsection*{Rozwiązanie}

Niech \(X\) będzie zmienną losową równą całkowitemu zyskowi graczy po milionie rozgrywek na maszynie, a \(X_i\) to zysk gracza przy \(i\)-tej rozgrywce. Na podstawie nierówności Chernoffa:

\begin{align*}
    P(X \geq 10000) &\leq \min_{t>0}\frac{E(e^{tX})}{e^{10000t}}\\
    &= \min_{t>0}\frac{E(e^{t\sum_{i=0}^{100000}X_i})}{e^{10000t}}\\
    &= \min_{t>0}\frac{E(\Pi_{i=0}^{100000}e^{X_i})}{e^{10000t}}\\
    &= \min_{t>0}\frac{E(e^{X_i})^{100000}}{e^{10000t}}
\end{align*}

Obliczymy najpierw funkcję tworzącą momenty pojedynczej rozgrywki na maszynie.

\[
E[e^{tX_i}] = \frac{1}{200}e^{99t} + \frac{4}{25}e^{2t} + \frac{167}{200}e^{-t}
\]

Wracając do poprzedniej nierówności:

\begin{align*}
    P(X \geq 10000) &\leq  \min_{t>0}\frac{E(e^{X_i})^{100000}}{e^{10000t}}\\
    &= \min_{t>0}\frac{(\frac{1}{200}e^{99t} + \frac{4}{25}e^{2t} + \frac{167}{200}e^{-t})^{100000}}{e^{10000t}}
\end{align*}

To nam daje pewne ograniczenie na \(t\). Można teraz wyznaczyć minimum numerycznie właściwie dowolną metodą - na przykład wykorzystać algorytm podobny do tego z zadania Cyklon z MP. Alternatywnie, tutaj jest \href{https://www.desmos.com/calculator/owyz2jam2l}{wykres}.

W ten sposób możemy oszacować minimalne ograniczenie górne dla \(t \sim 0.000577932\), co daje nam \(P(X \geq 10000) \leq \sim0.00015857\).

\subsection*{Odpowiedź}
\[
\boxed{P(X \geq 10000) \leq 0.00015857 \pm 0.00000001}
\]
